%% 第 19 章：自伴算子
%% 本文件由 tools/md2tex.py 自动生成，请勿直接编辑。

\chapter{自伴算子}


\begin{jdefn}{自伴算子的定义}{r:20:1}
若：

\[
T=T^*,
\]

则称 $T$ 是自伴算子。

在规范正交基下，自伴算子的矩阵满足：

\[
A=A^*.
\]

实数情形中，这就是对称矩阵；复数情形中，这就是 Hermitian 矩阵（共轭对称阵）。
\end{jdefn}

\begin{jprop}{自伴算子的基本性质}{r:20:2}
若 $S,T$ 都是自伴算子，则：

\[
S+T
\]

也是自伴算子。

若 $a\in\mathbb R$，则：

\[
aT
\]

也是自伴算子。

两个自伴算子的乘积不一定自伴。若额外满足：

\[
ST=TS,
\]

则：

\[
(ST)^*=T^*S^*=TS=ST,
\]

所以 $ST$ 自伴。
\end{jprop}

\begin{jprop}{自伴算子的特征值是实数}{r:20:3}
若：

\[
Tv=\lambda v,
\qquad
v\ne0,
\]

且 $T$ 自伴，则：

\[
\lambda\in\mathbb R.
\]
\end{jprop}

\begin{jproof}{证明思路}
计算：

\[
\langle Tv,v\rangle.
\]

一方面，由特征值关系，它等于 $\lambda\langle v,v\rangle$。

另一方面，由自伴性：

\[
\langle Tv,v\rangle
=
\langle v,Tv\rangle,
\]

所以它等于自己的共轭。

由于：

\[
\langle v,v\rangle>0,
\]

只能有：

\[
\lambda=\overline{\lambda}.
\]

因此 $\lambda$ 是实数。
\end{jproof}

\begin{jprop}{复内积空间中自伴的刻画}{r:20:4}
在复内积空间中：

\[
T\text{ 自伴}
\iff
\langle Tv,v\rangle\in\mathbb R
\text{ 对所有 }v\in V.
\]
\end{jprop}

\begin{jproof}{证明思路}
自伴时，内积可以直接看出等于自己的共轭，因此为实数。

反过来，若所有：

\[
\langle Tv,v\rangle
\]

都为实数，则分别对 $u+v$ 与 $u+iv$ 使用该条件。

展开后可分别得到：

\[
\langle Tu,v\rangle
=
\langle u,Tv\rangle.
\]

这正是：

\[
T=T^*.
\]
\end{jproof}

\begin{jprop}{复内积空间中的零二次型结论}{r:20:5}
若 $V$ 是复内积空间，且：

\[
\langle Tv,v\rangle=0
\]

对所有 $v\in V$ 成立，则：

\[
T=0.
\]
\end{jprop}

\begin{jproof}{证明思路}
分别代入 $u+v$ 与 $u+iv$。

前者表明：

\[
\langle Tu,v\rangle
\]

的实部为零；后者表明其虚部也为零。

因此：

\[
\langle Tu,v\rangle=0
\]

对所有 $u,v$ 成立。

特别取：

\[
v=Tu,
\]

得到：

\[
\|Tu\|^2=0.
\]

所以 $Tu=0$ 对所有 $u$ 成立，即：

\[
T=0.
\]

实内积空间中此结论不成立。平面上的 $90^\circ$ 旋转算子非零，但对每个 $v$ 都满足：

\[
\langle Tv,v\rangle=0.
\]
\end{jproof}
